% Generated by latex/build-latex.mjs from data/cv.json — edit the JSON, not this file.
% Compile with XeLaTeX (on Overleaf: Menu → Compiler → XeLaTeX), twice: the
% line joining roles at the same company is placed on the second run.
\documentclass[10pt,a4paper,sans]{moderncv}
\moderncvstyle{classic}
\moderncvcolor{black}
\usepackage[scale={0.8,0.87}]{geometry}  % a little less top/bottom margin, so sections fit where they belong
\usepackage{fontspec}
\setsansfont{IBMPlexSans-Regular.otf}[BoldFont=IBMPlexSans-SemiBold.otf, ItalicFont=IBMPlexSans-Italic.otf, BoldItalicFont=IBMPlexSans-SemiBoldItalic.otf]
\setmainfont{IBMPlexSans-Regular.otf}[BoldFont=IBMPlexSans-SemiBold.otf, ItalicFont=IBMPlexSans-Italic.otf, BoldItalicFont=IBMPlexSans-SemiBoldItalic.otf]
\setmonofont{cmuntt.otf}
\usepackage{polyglossia}
\setmainlanguage{greek}
\setotherlanguage{english}
\newfontfamily\greekfont{IBMPlexSans-Regular.otf}[BoldFont=IBMPlexSans-SemiBold.otf, ItalicFont=IBMPlexSans-Italic.otf, BoldItalicFont=IBMPlexSans-SemiBoldItalic.otf]
\newfontfamily\greekfontsf{IBMPlexSans-Regular.otf}[BoldFont=IBMPlexSans-SemiBold.otf, ItalicFont=IBMPlexSans-Italic.otf, BoldItalicFont=IBMPlexSans-SemiBoldItalic.otf]
\newfontfamily\greekfonttt{cmuntt.otf}
\newfontfamily\plexlight{IBMPlexSans-Light.otf}
\usepackage{array}
\usepackage{qrcode}
\usepackage{tikz}
\usepackage{needspace}
% Never leave a section heading alone at the foot of a page.
\let\cvsectionplain\section
\renewcommand*{\section}[1]{\needspace{5\baselineskip}\cvsectionplain{#1}}
\setlength{\hintscolumnwidth}{3.1cm}
\renewcommand*{\firstnamestyle}[1]{{\fontsize{32}{36}\plexlight\textcolor{firstnamecolor}{#1}}}
\renewcommand*{\lastnamestyle}[1]{{\fontsize{32}{36}\mdseries\textcolor{lastnamecolor}{#1}}}
\renewcommand*{\titlefont}{\large\mdseries\itshape}

% Marks and a hairline joining consecutive roles at one company.
\newcommand*{\cvmark}[1]{\tikz[remember picture,overlay]\coordinate (#1);}
\newcommand*{\cvlink}[2]{\tikz[remember picture,overlay]\draw[black!28, line width=0.5pt]
  ([xshift=1.6pt,yshift=-4pt]#1) -- ([xshift=1.6pt,yshift=10pt]#2);}

% Header: name and headline on the left; contact details on the right as a
% two-column table (icon, text), so every line starts at the same place.
\makeatletter
\renewcommand*{\makecvhead}{%
  \recomputecvlengths%
  \@initializebox{\makecvheaddetailsbox}%
  \savebox{\makecvheaddetailsbox}{{\small\upshape\renewcommand{\arraystretch}{1.3}%
    \begin{tabular}[b]{@{}l@{\hspace{0.65em}}>{\color{black!80}}l@{}}
      \makebox[1.1em][c]{\color{black!45}\faIcon{location-dot}} & Πάτρα, Αχαΐα, Ελλάδα\\
      \makebox[1.1em][c]{\color{black!45}\faIcon{mobile-screen}} & \href{tel:+306978036458}{+30 697 803 6458}\\
      \makebox[1.1em][c]{\color{black!45}\faIcon[regular]{envelope}} & \href{mailto:contact@ialexopoulos.org}{contact@ialexopoulos.org}\\
      \makebox[1.1em][c]{\color{black!45}\faIcon{earth-europe}} & \href{https://www.ialexopoulos.org/el/}{ialexopoulos.org}\\
      \makebox[1.1em][c]{\color{black!45}\faIcon{linkedin-in}} & \href{https://linkedin.com/in/ialexop}{linkedin.com/in/ialexop}\\
      \makebox[1.1em][c]{\color{black!45}\faIcon{github}} & \href{https://github.com/johnnypatras}{github.com/johnnypatras}\\
    \end{tabular}}}%
  \@initializelength{\makecvheaddetailswidth}%
  \settowidth{\makecvheaddetailswidth}{\usebox{\makecvheaddetailsbox}}%
  \setlength{\makecvheadnamewidth}{\textwidth-\makecvheaddetailswidth-1em}%
  \@initializebox{\makecvheadnamebox}%
  \savebox{\makecvheadnamebox}{%
    \begin{minipage}[b]{\makecvheadnamewidth}%
      \raggedright\firstnamestyle{\@firstname\ }\lastnamestyle{\@lastname}%
      \ifthenelse{\equal{\@title}{}}{}{\\[1em]\titlestyle{\@title}}%
    \end{minipage}}%
  \usebox{\makecvheadnamebox}\hfill\usebox{\makecvheaddetailsbox}%
  \par\vspace{2.5em}}
\makeatother

\name{Ιωάννης}{Αλεξόπουλος}
\title{Διοίκηση λειτουργιών \& επιχειρηματική ανάπτυξη\texorpdfstring{\\}{ · }Μηχανικός Η/Υ και Πληροφορικής}  % role, then degree on its own line

\begin{document}
% moderncv loads hyperref at \begin{document}, so the PDF metadata is set here.
\hypersetup{pdftitle={Ιωάννης Αλεξόπουλος -- Βιογραφικό σημείωμα}, pdfauthor={Ιωάννης Αλεξόπουλος}, pdflang={el}}
\makecvtitle

\section{Επαγγελματικό προφίλ}
Περισσότερα από 15 χρόνια εμπειρίας στη διοίκηση ομάδων, στην οργάνωση της λειτουργίας και στην ανάπτυξη πωλήσεων. Στη NOBACCO, εξέλιξη από τη θέση του υπευθύνου καταστήματος σε αυτή του περιφερειακού διευθυντή, με ευθύνη για τα καταστήματα της περιοχής, ομάδα 30 εργαζομένων, οκτώ καταστήματα franchise και τη χονδρική διάθεση σε Πελοπόννησο, Δυτική Ελλάδα και Ιόνια Νησιά. Δίπλωμα Μηχανικού Η/Υ και Πληροφορικής του Πανεπιστημίου Πατρών· συνδυασμός τεχνικού υποβάθρου και εμπειρίας στην επιχειρηματική ανάπτυξη. Παρακολούθηση προγραμμάτων πιστοποίησης στη διαχείριση έργων και στις σύγχρονες τεχνολογίες, με ενδιαφέρον για θέσεις που συνδυάζουν τη διοίκηση λειτουργιών με την τεχνολογία.

\section{Επαγγελματική εμπειρία}
  \cvitem{}{\textbf{NOBACCO} \textperiodcentered{} 2015–2025}
  \cventry{\cvmark{g0r0}2019–2025}{Περιφερειακός Διευθυντής \& Υπεύθυνος Επιχειρηματικής Ανάπτυξης}{}{}{}{
    \begin{itemize}
      \item Εποπτεία της λειτουργίας των καταστημάτων της περιοχής: τήρηση διαδικασιών, λειτουργικό κόστος, απόδοση (ROI)
      \item Διοίκηση ομάδας 30 εργαζομένων — εκπαίδευση, ετήσιες αξιολογήσεις, προσλήψεις
      \item Έναρξη λειτουργίας οκτώ καταστημάτων franchise και υποστήριξή τους με εκπαίδευση, καθοδήγηση και παρακολούθηση απόδοσης
      \item Συμφωνίες χονδρικής διάθεσης των προϊόντων NOBACCO με χονδρεμπόρους σε Αχαΐα, Αιτωλοακαρνανία, Αργολίδα, Κορινθία, Λακωνία, Μεσσηνία και Ιόνια Νησιά
    \end{itemize}}
  \cventry{\cvmark{g0r1}\cvlink{g0r0}{g0r1}2015–2019}{Υπεύθυνος Καταστήματος}{}{}{}{
    \begin{itemize}
      \item Επίτευξη των μηνιαίων, τριμηνιαίων και ετήσιων στόχων του καταστήματος
      \item Διαχείριση προσωπικού, ωραρίων και αποθήκης
      \item Τήρηση των προτύπων εξυπηρέτησης και της εταιρικής κουλτούρας· μεταφορά της εταιρικής ενημέρωσης στην ομάδα
    \end{itemize}}
  \cventry{2010–2012}{Διαχειριστής Συστημάτων}{Internet café \textperiodcentered{} Κάτω Αχαΐα}{}{}{
    \begin{itemize}
      \item Συντήρηση και αναβάθμιση υπολογιστών, περιφερειακών και τοπικού δικτύου
      \item Τακτική λήψη αντιγράφων ασφαλείας και προστασία δεδομένων
      \item Τεχνική υποστήριξη και άμεση επίλυση προβλημάτων
    \end{itemize}}
  \cventry{2004–2007}{Εκπρόσωπος Εξυπηρέτησης Πελατών}{ΟΤΕ Α.Ε.}{}{}{
    \begin{itemize}
      \item Διαχείριση αιτημάτων, παραπόνων και κλιμακώσεων πελατών
      \item Ενημέρωση πελατών για προϊόντα και υπηρεσίες
    \end{itemize}}

\section{Δεξιότητες}
  \cvitem{Σε εξέλιξη}{Professional Scrum Master (PSM) \textperiodcentered{} CS50P και CS50AI, Harvard \textperiodcentered{} Oracle Java Foundations}
  \cvitem{Διοίκηση}{Στρατηγική ηγεσία και ανάπτυξη ομάδων \textperiodcentered{} Επιχειρηματικός σχεδιασμός και διαχείριση συνεργασιών \textperiodcentered{} Αναλυτική επίλυση προβλημάτων και βελτίωση διαδικασιών}
  \cvitem{Εφαρμοσμένη εμπειρία}{Python (έξι μήνες σε εταιρικό έργο)}
  \cvitem{Ακαδημαϊκό επίπεδο}{Java \textperiodcentered{} JavaScript \textperiodcentered{} SQL \textperiodcentered{} MATLAB \textperiodcentered{} C \textperiodcentered{} C++ \textperiodcentered{} HTML5 \textperiodcentered{} CSS}
  \cvitem{Εργαλεία}{Linux (Debian, Fedora, Ubuntu) \textperiodcentered{} Git \& GitHub \textperiodcentered{} VS Code \textperiodcentered{} LaTeX \textperiodcentered{} Microsoft Office}

\section{Σπουδές}
  \cventry{2019}{Δίπλωμα Μηχανικού Η/Υ και Πληροφορικής}{Πανεπιστήμιο Πατρών \textperiodcentered{} Πολυτεχνική Σχολή}{}{}{Διπλωματική εργασία: «Low Code Development Platforms»}

\section{Σεμινάρια}
  \cvitem{2017–2024}{Εξυπηρέτηση πελατών – Ελληνικό Ινστιτούτο Εξυπηρέτησης Πελατών \textperiodcentered{} 2017–2019 και 2022–2024}
  \cvitem{2018}{Στρατηγική και εξειδικευμένες τεχνικές πωλήσεων – NOBACCO Α.Ε. / BAT Hellas}
  \cvitem{2017}{Βέλτιστες πρακτικές εξυπηρέτησης πελατών – NOBACCO Α.Ε.}

\section{Ξένες γλώσσες}
  \cvitem{Ελληνικά}{Μητρική}
  \cvitem{Αγγλικά}{C2 \textperiodcentered{} Άριστη γνώση \textperiodcentered{} Certificate of Proficiency in English, University of Michigan}
  \cvitem{Γερμανικά}{C1 \textperiodcentered{} Πολύ καλή γνώση \textperiodcentered{} Zentrale Mittelstufenprüfung, Goethe-Institut}

% QR code to the matching language page of the website, for printed copies:
% the code sits in the date column, its caption centred beside it.
\bigskip
\noindent\parbox[c]{\hintscolumnwidth}{\raggedleft\qrcode[height=1.9cm]{https://www.ialexopoulos.org/el/}}%
\hspace{\separatorcolumnwidth}%
\parbox[c]{\maincolumnwidth}{\small\color{black!70}Ηλεκτρονική έκδοση του βιογραφικού, με αναλυτική περιγραφή των αρμοδιοτήτων\\[2pt]{\color{black}\href{https://www.ialexopoulos.org/el/}{ialexopoulos.org/el/}}}

\end{document}
